% Automated Validation & Metric Collection System
% Paper-ready Methods and Implementation report
\documentclass[11pt]{article}
\usepackage{graphicx}
\usepackage{booktabs}
\usepackage{amsmath}
\usepackage{hyperref}
\usepackage{caption}
\usepackage{geometry}
\geometry{margin=1in}

\title{Automated Validation and Metric Collection System (AVMCS): \newline Methods and Reference Implementation}
\author{(Repository: drakeDjkw/Automated-Validation-Metric-Collection-System)}
\date{\today}

\begin{document}
\maketitle

\begin{abstract}
We present a compact, modular framework for automated data validation, batched model evaluation, metric collection, and reproducible experiment logging. The included reference implementation (Python, numpy/pandas) demonstrates core components that are essential for trustworthy ML system development: deterministic data checks, per-batch evaluation, metric persistence, and CI-driven verification. We provide a detailed Methods description, a worked synthetic experiment, and reproducible artifacts (CSV, PNG) suitable for inclusion in academic manuscripts.
\end{abstract}

\section{Introduction}
Modern machine learning systems require continuous validation of data and automatic monitoring of model performance to avoid silent degradation in production. AVMCS provides a minimal, extensible skeleton to perform: (1) data validation gates, (2) batched inference and metric computation, and (3) simple artifact logging for downstream analysis. Its design emphasizes reproducibility, testability, and easy integration with experiment tracking and dashboard tools.

\section{Design and Architecture}
The system implements a linear evaluation pipeline: \textbf{Data source} $\rightarrow$ \textbf{Validation layer} $\rightarrow$ \textbf{Model (inference)} $\rightarrow$ \textbf{Evaluation (metrics)} $\rightarrow$ \textbf{Metrics store} $\rightarrow$ \textbf{Dashboard / Alerts}. Critical errors are logged and can halt downstream processing.

Core code-to-concept mapping (file references are relative to the repository root):
\begin{itemize}
  \item \texttt{src/avmcs/validation.py} -- data checks (missing values, type checks, range checks).
  \item \texttt{src/avmcs/metrics.py} -- implementation of MAE and RMSE.
  \item \texttt{src/avmcs/pipeline.py} -- batched evaluation logic and \texttt{DummyModel} demonstration.
  \item \texttt{src/avmcs/storage.py} -- JSON storage for metrics.
  \item \texttt{scripts/run_pipeline.py} -- runnable demo pipeline.
  \item \texttt{scripts/run_experiments.py} -- synthetic experiments producing \texttt{experiments/results.csv} and \texttt{experiments/metrics.png}.
  \item \texttt{tests/test_pipeline.py} -- unit test ensuring correctness of the evaluation pipeline.
\end{itemize}

\section{Methods and Implementation}
This section maps algorithms and equations to repository code.

\subsection{Data validation}
The validation module implements deterministic rules applied to tabular datasets (pandas.DataFrame). Each check is specified per column via a small dictionary of expectations (type and allowable range). The primary function is:

\begin{verbatim}
def validate_data(df: pd.DataFrame, checks: Dict[str, Any] = None) -> Dict[str, Any]
\end{verbatim}

Where `checks` follows the shape:

\begin{verbatim}
{
    "temperature": {"type": "numeric", "range": [-50, 60]},
    "humidity": {"type": "numeric", "range": [0, 100]},
}
\end{verbatim}

The function performs:
\begin{itemize}
  \item Missing-value detection: per-column counts (returns `missing_values`).
  \item Type checks: currently supports `numeric` via `pandas.api.types.is_numeric_dtype` (returns `type_errors`).
  \item Range checks: reports `range_errors` when any value falls outside `[min,max]`.
\end{itemize}

Return format: a dictionary of error categories, for example:

\begin{verbatim}
{
  "missing_values": {"colA": 5},
  "range_errors": {"temperature": "Values outside [-50, 60]"}
}
\end{verbatim}

Design choice: simple, deterministic checks that are fast and interpretable. In production, these checks can be extended to include schema validation, categorical domain checks, and statistical profiling.

\subsection{Metrics}
We use two core regression metrics implemented in \texttt{src/avmcs/metrics.py}: MAE and RMSE. For a prediction set $\{\hat{y}_i\}_{i=1}^N$ and ground truth $\{y_i\}_{i=1}^N$:
\begin{align}
\text{MAE} &= \frac{1}{N} \sum_{i=1}^N |y_i - \hat{y}_i|, \\
\text{RMSE} &= \sqrt{\frac{1}{N} \sum_{i=1}^N (y_i - \hat{y}_i)^2 }.
\end{align}

The implementation verifies shape compatibility and returns a small JSON-serializable dict: \texttt{{"MAE": float, "RMSE": float}}.
In the code the function signature is:

\begin{verbatim}
def compute_metrics(y_true: np.ndarray, y_pred: np.ndarray) -> Dict[str, float]
\end{verbatim}

Important implementation notes:
\begin{itemize}
  \item `numpy.asarray` is used to convert inputs; the function raises `ValueError` if shapes differ.
  \item The function returns Python floats (safe for JSON serialization) and is intentionally minimal so users can add additional metrics (SSIM, PSNR, FID) as needed for images or generative models.
\end{itemize}

\subsection{Batched evaluation pipeline}
The evaluation pipeline is implemented in \texttt{run\_pipeline(model, X, y, batch\_size)}. For each batch of inputs it obtains predictions and computes per-batch metrics. It aggregates metrics by averaging across batches, returning both per-batch metrics and aggregated metrics. In the code the signature is:

\begin{verbatim}
def run_pipeline(model: Callable[..., Any], X: np.ndarray, y: np.ndarray, batch_size: int = 32) -> Dict[str, Any]
\end{verbatim}

Pseudocode follows:

\begin{verbatim}
for i in range(0, N, batch_size):
    xb = X[i:i+batch_size]
    yb = y[i:i+batch_size]
    y_pred = model.predict(xb)  # or model(xb)
    metrics_batch = compute_metrics(yb, y_pred)
    append metrics_batch
aggregated = mean_over_batches(batch_metrics)
\end{verbatim}

Note: for production-grade reporting we recommend aggregating across all predictions (global MAE/RMSE) rather than averaging per-batch statistics; the repository uses the simpler mean-of-batches for clarity.

Additional pipeline details:
\begin{itemize}
  \item `model` may be a callable (function) returning predictions or an object exposing a `.predict(X)` method; the implementation checks `hasattr(model, 'predict')` to handle both.
  \item The function returns a dict with keys `batch_metrics` (list of per-batch metric dicts) and `aggregated` (dict of averaged metrics).
  \item Example aggregated output (JSON):
  \begin{verbatim}
  {
    "batch_metrics": [{"MAE": 1.23, "RMSE": 1.56}, ...],
    "aggregated": {"MAE": 1.25, "RMSE": 1.60}
  }
  \end{verbatim}
\end{itemize}

\subsection{Storage and reproducibility}
Metrics are saved as JSON via \texttt{save\_metrics}. The experiments script writes a CSV summary and a PNG figure to the \texttt{experiments/} directory. The repository also includes a \texttt{requirements.txt} and a CI workflow (GitHub Actions) that runs linting, tests, and a demo run to ensure baseline reproducibility on every push.

The storage API signatures:
\begin{verbatim}
def save_metrics(metrics: Any, filename: str = "results.json") -> None
def load_metrics(filename: str = "results.json") -> Any
\end{verbatim}

The JSON file is intentionally schema-light; for research-grade tracking we recommend adopting MLflow, W&B, or saving a richer schema with timestamps, dataset/model hashes, and hyperparameters.

\section{Experimental Protocol (synthetic)}
We include a synthetic experiment to illustrate the pipeline and to produce figures/tables for the manuscript. Dataset: features sampled from an isotropic Gaussian. Target: sum of features plus Gaussian noise. Model: a trivial \texttt{DummyModel} that returns average-of-features plus an additive \emph{bias} parameter. We vary the bias to see how systematic model offset affects MAE and RMSE.

Protocol:
\begin{itemize}
  \item For each bias value in a grid (e.g., -1.0 to 1.0), and for several random seeds:
  \item Generate dataset with fixed sample size and feature count.
  \item Compute predictions via \texttt{DummyModel(bias)} and evaluate with \texttt{run\_pipeline}.
  \item Save aggregated metrics per run to \texttt{experiments/results.csv}.
  \item Compute mean and standard deviation across seeds for each bias and plot MAE/RMSE vs bias.
\end{itemize}

The figure and the CSV are committed in the repository (\texttt{experiments/metrics.png}, \texttt{experiments/results.csv}). The README includes the computed table of mean $\pm$ std values.

Concrete experiment parameters used in \texttt{scripts/run\_experiments.py}:
\begin{itemize}
  \item Number of samples per run: 500
  \item Number of features: 5
  \item Noise std for target: 0.5
  \item Bias grid: 9 values from -1.0 to 1.0 (inclusive)
  \item Seeds per bias: 3 (seeds 0,1,2)
  \item Batch size used in pipeline: 64
\end{itemize}

These choices are kept small to make the example fast and reproducible in CI; you can increase `n` and `seeds` for tighter estimates.

\begin{figure}[ht]
  \centering
  \includegraphics[width=0.7\linewidth]{experiments/metrics.png}
  \caption{MAE and RMSE vs model bias on the synthetic dataset. Points reflect averages across random seeds.}
  \label{fig:metrics}
\end{figure}

\section{Results (summary)}
Representative results from the synthetic experiments (mean $\pm$ std across seeds) are summarized in the README and reproduced here: Table~\ref{tab:summary}.

\begin{table}[ht]
\centering
\begin{tabular}{rcc}
\toprule
Bias & MAE (mean $\pm$ std) & RMSE (mean $\pm$ std)\\
\midrule
-1.00 & 1.6410 $\pm$ 0.0409 & 2.0513 $\pm$ 0.0392\\
-0.75 & 1.5488 $\pm$ 0.0292 & 1.9475 $\pm$ 0.0310\\
-0.50 & 1.4830 $\pm$ 0.0207 & 1.8714 $\pm$ 0.0252\\
-0.25 & 1.4433 $\pm$ 0.0127 & 1.8264 $\pm$ 0.0258\\
0.00 & 1.4344 $\pm$ 0.0041 & 1.8148 $\pm$ 0.0333\\
0.25 & 1.4563 $\pm$ 0.0180 & 1.8372 $\pm$ 0.0438\\
0.50 & 1.5075 $\pm$ 0.0340 & 1.8924 $\pm$ 0.0545\\
0.75 & 1.5863 $\pm$ 0.0466 & 1.9779 $\pm$ 0.0641\\
1.00 & 1.6883 $\pm$ 0.0584 & 2.0899 $\pm$ 0.0722\\
\bottomrule
\end{tabular}
\caption{MAE and RMSE (mean $\pm$ std) across seeds for each model bias (synthetic experiment).}
\label{tab:summary}
\end{table}

\section{Discussion and Extensions}
The repository demonstrates a minimal but practical approach to automated validation and metric collection. For production or domain-specific research, recommended extensions include:
\begin{itemize}
  \item Statistical drift detection (KL/JS divergence, KS tests, PSI) to detect distributional shifts between training and incoming data.
  \item Integration with experiment tracking systems (MLflow \cite{mlflow}, Weights & Biases \cite{wandb}) for versioned metrics, model artifacts, and hyperparameters.
  \item More advanced validation with expectation frameworks (e.g., Great Expectations) to encode and test complex dataset invariants.
  \item Domain-specific metrics: for image forecasting or generative models, add SSIM \cite{ssim}, PSNR, FID \cite{fid}, or task-specific scores.
\end{itemize}

\section{Reproducibility and Availability}
All code and artifacts are available in the repository: \url{https://github.com/drakeDjkw/Automated-Validation-Metric-Collection-System}. Reproducible commands and required dependencies are documented in \texttt{README.md}. The CI workflow ensures that linting and tests run on each push and PR.

\section{References}
\begin{thebibliography}{9}
\bibitem{goodfellow2014gan}
Ian J. Goodfellow, Jean Pouget-Abadie, Mehdi Mirza, Bing Xu, David Warde-Farley, Sherjil Ozair, Aaron Courville, and Yoshua Bengio. "Generative Adversarial Nets." Advances in Neural Information Processing Systems 27 (NIPS 2014).

\bibitem{shi2015convlstm}
Xingjian Shi, Zhourong Chen, Hao Wang, Dit-Yan Yeung, Wai-kin Wong, and Wang-chun Woo. "Convolutional LSTM Network: A Machine Learning Approach for Precipitation Nowcasting." Advances in Neural Information Processing Systems 2015.

\bibitem{wang2004ssim}
Z. Wang, A. C. Bovik, H. R. Sheikh, and E. P. Simoncelli. "Image quality assessment: from error visibility to structural similarity." IEEE Transactions on Image Processing, 2004.

\bibitem{fid}
Martin Heusel, Hubert Ramsauer, Thomas Unterthiner, Bernhard Nessler, and Sepp Hochreiter. "GANs Trained by a Two Time-Scale Update Rule Converge to a Local Nash Equilibrium." In Advances in Neural Information Processing Systems, 2017.

\bibitem{sculley2015hidden}
D. Sculley, Gary Holt, Daniel Golovin, Eugene Davydov, Todd Phillips, Dietmar Ebner, Vinay Chaudhary, Michael Young, Jean-François Crespo, and Dan Dennison. "Hidden Technical Debt in Machine Learning Systems." In NIPS Workshop on ML and Systems, 2015.

\bibitem{mlflow}
MLflow. https://mlflow.org (software). Accessed 2026.

\bibitem{wandb}
Weights \& Biases. https://wandb.ai (software). Accessed 2026.

\bibitem{greatexp}
Great Expectations. https://greatexpectations.io (software). Accessed 2026.
\end{thebibliography}

\end{document}
